\begin{frame}{Struttura della Presentazione}
    \begin{enumerate}
        \item \textbf{La corrispondenza di Curry-Howard (CH)}
        \begin{itemize}
            \item Il ponte strutturale tra logica intuizionista e teoria dei tipi.
            \item Isomorfismo statico (formule $\leftrightarrow$ tipi) e dinamico (cut-elimination $\leftrightarrow$ $\beta$-riduzione).
        \end{itemize}
        
        \item \textbf{Che cos'è Lean 4 e cosa ha a che vedere con CH}
        \begin{itemize}
            \item Il Calculus of Inductive Constructions (CiC) e la gerarchia degli universi (\texttt{Prop} vs \texttt{Type}).
            \item Small Trusted Kernel (principio di de Bruijn) e garanzie della Trusted Computing Base (TCB).
        \end{itemize}
        
        % TODO: modificare titolo
        \item \textbf{Cosa si può effettivamente provare di CH in Lean (Capitolo 8)}
        \begin{itemize}
            \item Deep embedding e gestione delle variabili mediante indici di de Bruijn.
            \item Morfismi costruttivi ed eseguibili \texttt{extract} ed \texttt{embed}.
            \item Metateoremi provati: involuzione sintattica, coerenza geometrica e commutatività dinamica.
        \end{itemize}
    \end{enumerate}
\end{frame}